\documentclass{article}

\usepackage[margin=1in]{geometry}
\usepackage{graphicx, amsmath, amsthm, amssymb, enumitem, lipsum}

\setlist[enumerate]{leftmargin=1.25cm}

\newcommand{\N}{\mathbb{N}}
\newcommand{\Z}{\mathbb{Z}}
\newcommand{\Q}{\mathbb{Q}}
\newcommand{\R}{\mathbb{R}}
\newcommand{\C}{\mathbb{C}}

\newcommand{\rotvert}{\rotatebox[origin=c]{90}{$\vert$}}

\newenvironment{note}[2][Note]{\begin{trivlist}
\item[\hskip \labelsep {\bfseries #1}\hskip \labelsep {\bfseries #2.}]}{\end{trivlist}}

\newenvironment{fact}[2][Fact]{\begin{trivlist}
\item[\hskip \labelsep {\bfseries #1}\hskip \labelsep {\bfseries #2.}]}{\end{trivlist}}

\newenvironment{definition}[2][Definition]{\begin{trivlist}
\item[\hskip \labelsep {\bfseries #1}\hskip \labelsep {\bfseries #2.}]}{\end{trivlist}}

\newenvironment{proposition}[2][Proposition]{\begin{trivlist}
\item[\hskip \labelsep {\bfseries #1}\hskip \labelsep {\bfseries #2.}]}{\end{trivlist}}

\newenvironment{principle}[2][Principle]{\begin{trivlist}
\item[\hskip \labelsep {\bfseries #1}\hskip \labelsep {\bfseries #2.}]}{\end{trivlist}}

\newenvironment{lemma}[2][Lemma]{\begin{trivlist}
\item[\hskip \labelsep {\bfseries #1}\hskip \labelsep {\bfseries #2.}]}{\end{trivlist}}

\newenvironment{theorem}[2][Theorem]{\begin{trivlist}
\item[\hskip \labelsep {\bfseries #1}\hskip \labelsep {\bfseries #2.}]}{\end{trivlist}}

\newenvironment{corollary}[2][Corollary]{\begin{trivlist}
\item[\hskip \labelsep {\bfseries #1}\hskip \labelsep {\bfseries #2.}]}{\end{trivlist}}

\newenvironment{envsection}[1]{\begin{trivlist}
\item[\hskip \labelsep {\bfseries #1}]}{\end{trivlist}}


\begin{document}
\setlength{\abovedisplayskip}{4pt}
\setlength{\belowdisplayskip}{4pt}


\begin{center}
    \Large{\textbf{Homework 2}}\\
    \normalsize February 5 - 12
\end{center}
\vspace{10pt}

\subsection*{Reading}

\begin{envsection}{1. CS229 Linear Algebra Reference:}
    Sections 1 and 2: Basic Concepts and Notation and Matrix Multiplication. Use this as a reference as you work on the homework problems. Section 2 does a great job outlining and explaining more about the usefulness of the language of linear algebra/matrix multiplication.
\end{envsection}

\noindent Work on the problems first, but if you finish:

\begin{envsection}{2. Rogers and Girolami:}
    Sections 1.4 - 1.6.
\end{envsection}

\begin{envsection}{3. CS229 Notes:}
    Read the introductions to Part I and Chapter 1 and Section 1.1: LMS algorithm.
\end{envsection}


% \begin{enumerate}
%     \item[\textbf{1.}] \textbf{CS229 Linear Algebra Reference:} Sections 1 and 2: Basic Concepts and Notation and Matrix Multiplication
%     \item[\textbf{2.}] 
% \end{enumerate}

\subsection*{Problems}

Finish Problems 1, 2, and 3 from Homework 1. For Problem 3, I'd like you to redo part (a) yourself even though we finished it together last meeting. Be prepared to present your solutions next meeting, if you can.


\end{document}